\chapter{Personal Risk Management and Insurance Strategy}

A financial plan assumes things go as expected. When they do not---a medical bill, a broken motorcycle, a missed freelance payment---the plan needs a way to absorb the hit without wiping out months of savings. Risk management here works at two levels: a cash buffer for ordinary disruptions, and insurance for larger, less predictable events.

\section{Emergency Fund Plan}

My current risk coverage across four areas: health is covered by the mandatory student Bảo hiểm Y tế (BHYT) for public hospitals plus a Bảo Việt An Gia Silver private plan (private hospitals, direct billing); income disruption risk is managed through three concurrent sources (no single source is more than 42\% of total income); the Datbike Quantum S2 has vehicle accident insurance; the family apartment has no rental exposure on my side.

The remaining gap is cash. Emergency fund target: \textbf{30,000,000~VND} (approximately 6 months of expenses at the 4,920,000~VND/month target budget), held entirely in cash or short-term bank deposits---not stocks or Bitcoin. Current cash-only cover is only 1.5 months (see Phase~3). Building this fund runs in parallel with the equipment upgrade goal.

\section{Insurance Coverage by Life Stage}

\begin{table}[h!]
\caption{Insurance Planning by Life Stage}
\label{tab:insurance_stages}
\begin{tabularx}{\linewidth}{@{}lLL@{}}
\toprule
\textbf{Life Stage} & \textbf{Insurance Coverage Needed} & \textbf{Rationale} \\
\midrule
Student (Ages 18--22)   & Mandatory student Bảo hiểm Y tế (BHYT) + Bảo Việt An Gia & Low-cost access to private clinics alongside public BHYT cover. \\[6pt]
First Job (Ages 23--27) & Employer-provided group health insurance (BCG benefit) & Full medical, dental, and outpatient cover funded by employer. \\[6pt]
Mid-Career (Ages 28+)   & Term Life + Critical Illness Insurance & Protects family dependents if the main earner dies or can no longer work. \\
\bottomrule
\end{tabularx}
\end{table}

At each stage, the type of risk that matters most changes. As a student with no dependents, the main risk is a large personal medical bill. At BCG, employer insurance covers this at no extra cost. From age 28, with a spouse, children, and aging parents depending on my income, the risk shifts to the family losing its main income source---which term life and critical illness insurance are designed to address.

\section{Evaluated Insurance Policy}

Selected: \textbf{Bảo Việt An Gia Silver Package}. Annual premium: 2,500,000 VND per year (approximately 208,000 VND per month). Coverage: inpatient and surgery up to 90,000,000 VND per year (surgery sub-limit: 45,000,000 VND); room and board up to 4,500,000 VND per day; direct billing at Hong Ngoc, Thu Cuc, and Vinmec hospitals in Hanoi. At 208,000 VND per month, this plan removes the need to pay upfront at private hospitals and covers most realistic medical events in Vietnam---making it one of the best-value decisions in this entire plan.

To see why this is worth 208,000 VND per month, consider what Bảo hiểm Y tế (BHYT) only coverage means in practice. BHYT sends patients to public hospitals, where outpatient queues commonly run two to four hours, shared wards are standard for inpatient stays, and access to specialist care requires an internal referral process that can add days to a diagnosis. For a student whose income depends on consistently showing up to an internship and delivering freelance projects on time, a two-day delay in getting a diagnosis is not just uncomfortable---it has a direct cost in lost work. The Bảo Việt An Gia plan allows same-day appointments at private hospitals with direct specialist access and no upfront payment. For 208,000 VND per month, it converts a potentially multi-day health disruption into a two-hour clinic visit.

Two practical renewal considerations: first, this plan should be reviewed annually. If the hospital network changes or new private health plans with better coverage-to-premium ratios become available, switching at renewal costs nothing. Second, when joining BCG after graduation in 2029, the employer group health insurance will replace this plan at zero personal cost. At that point, the personal Bảo Việt An Gia policy should be cancelled---there is no benefit in paying twice for overlapping coverage.

